
\documentclass[11pt,a4paper]{article}
\usepackage[utf8]{inputenc}
\usepackage[italian]{babel}
\usepackage{amsmath,amssymb}
\usepackage{geometry}
\geometry{margin=25mm}
\usepackage{xcolor}

\definecolor{primarycolor}{RGB}{41, 128, 185}
\definecolor{boxbg}{RGB}{245, 247, 250}

\title{\textbf{Metodo di Trachtenberg -- Appunti di Calcolo Mentale}}
\author{}
\date{}

\begin{document}
\maketitle

\section{Moltiplicazione per 11}

\textbf{Regola:} Si scrive l'ultima cifra del numero (quella delle unità) inalterata. Poi, per ogni cifra successiva, si somma la cifra stessa alla cifra alla sua destra. Infine, si riporta l'ultima cifra a sinistra (la prima del numero originario).

\textbf{Esempio:} Calcoliamo $432 \times 11$.
\begin{enumerate}
    \item Scriviamo l'unità finale: \textbf{2}.
    \item Sommiamo la cifra centrale con la destra ($3 + 2$): \textbf{5}.
    \item Sommiamo la prima cifra con la destra ($4 + 3$): \textbf{7}.
    \item Scriviamo la prima cifra iniziale: \textbf{4}.
\end{enumerate}
Risultato: $432 \times 11 = \mathbf{4752}$.

\section{Moltiplicazione per 12}

\textbf{Regola:} Si raddoppia ogni cifra e sommandola al suo vicino di destra (la cifra immediatamente successiva verso destra). 

\textbf{Esempio:} Calcoliamo $314 \times 12$.
\begin{enumerate}
    \item Per l'ultima cifra (4): $4 \times 2 + 0 = \mathbf{8}$.
    \item Per la cifra centrale (1): $1 \times 2 + 4 = \mathbf{6}$.
    \item Per la prima cifra (3): $3 \times 2 + 1 = \mathbf{7}$.
    \item Per il bordo a sinistra (aggiungendo uno zero iniziale implicito): $0 \times 2 + 3 = \mathbf{3}$.
\end{enumerate}
Risultato: $314 \times 12 = \mathbf{3768}$.

\section{Moltiplicazione per 6}

\textbf{Regola:} Si prende ogni cifra, si somma la metà della cifra alla sua destra (arrotondata per difetto se dispari) e si aggiunge 5 se la cifra stessa è un numero dispari.

\textbf{Esempio:} Calcoliamo $43 \times 6$.
\begin{enumerate}
    \item Cifra 3 (dispari): aggiungiamo 5 ($5$). A destra non c'è nulla (metà = 0). Totale: $5$.
    \item Cifra 4 (pari): aggiungiamo 0 ($0$). A destra c'è il 3, la cui metà intera è $1$. Sommiamo la cifra stessa ($4 + 1 = 5$).
    \item Per la cifra a sinistra dello zero iniziale: metà del 4 è $2$.
\end{enumerate}
Risultato: $43 \times 6 = \mathbf{258}$.

\end{document}